\section{Methodology}
\label{sec:method}

\subsection{Problem setup}
\label{sec:setup}

Let $T \in \mathbb{R}_{>0}$ be a failure time and $C \in \mathbb{R}_{>0}$ be
an independent right-censoring time, both depending on a covariate
$Z \in \mathbb{R}^{d}$. The observed data are i.i.d.\ tuples
\[
    \mathcal{D} = \{(X_i,\, \delta_i,\, Z_i)\}_{i=1}^{n},
    \qquad
    X_i = \min(T_i, C_i),
    \quad
    \delta_i = \mathbf{1}\{T_i \le C_i\}.
\]
We assume \emph{conditionally independent censoring}, $T \perp C \mid Z$
(the uninformative case); the shared-frailty informative-censoring extension
is sketched in Section~\ref{sec:discussion}. We jointly learn two generative
models for $T \mid Z$ and $C \mid Z$ such that both the conditional density
and conditional survival function admit closed forms (used in the censored
log-likelihood) and such that samples can be drawn with exact gradients (used
in the distributional regularizer). We meet both requirements with a real
conditional normalizing flow whose forward direction is closed-form and whose
inverse is reattached to gradients via implicit reparameterization. The two
flows are denoted $h_{\theta_1}$ for $T$ and $h_{\theta_2}$ for $C$ and have
disjoint parameters; for brevity we describe a single flow in
Section~\ref{sec:flow}--\ref{sec:samp} and write $\theta = \theta_1$ or
$\theta_2$.

\subsection{Conditional monotone normalizing flow}
\label{sec:flow}

\paragraph{Forward map.}
We parameterize the flow on the log-time axis with a $K$-component
tanh-of-affine basis on top of an MLP backbone. Given a covariate $z$, the
forward map $h_\theta(\,\cdot\,;\, z): \mathbb{R} \to \mathbb{R}$ is
\begin{equation}
    h_\theta(\log t;\, z)
    \;=\;
    a(z)\,\log t \;+\; b(z)
        \;+\; \sum_{k=1}^{K} \alpha_k(z)\, \tanh\!\left(\frac{\log t - \mu_k(z)}{s_k(z)}\right),
    \label{eq:flow}
\end{equation}
where $a, b, \{\alpha_k, \mu_k, s_k\}_{k=1}^{K}$ are produced by a single
shared MLP $\psi_\theta: \mathbb{R}^{d} \to \mathbb{R}^{2 + 3K}$ followed by
output-specific reparameterizations. In the released implementation $\psi$
has the architecture
\[
    z
    \;\xrightarrow{\;\;\text{Linear}(d \to 32)\;\;}
    \;\xrightarrow{\;\;\tanh\;\;}
    \;\xrightarrow{\;\;\text{Linear}(32 \to 32)\;\;}
    \;\xrightarrow{\;\;\tanh\;\;}
    \;\xrightarrow{\;\;\text{Linear}(32 \to 2 + 3K)\;\;}
    \;\;,
\]
and we split its $2 + 3K$ outputs into raw heads
$(\widetilde{a}, b, \widetilde{\alpha}_{1:K}, \mu_{1:K}, \widetilde{s}_{1:K})$.
The constraint reparameterizations are
\begin{equation}
    a(z) \;=\; \operatorname{softplus}(\widetilde{a}(z)) + a_0,
    \quad
    \alpha_k(z) \;=\; \alpha_{\max}\,\operatorname{softplus}(\widetilde{\alpha}_k(z)),
    \quad
    s_k(z) \;=\; \operatorname{softplus}(\widetilde{s}_k(z)) + s_0,
    \label{eq:reparam}
\end{equation}
with shifts $a_0 = 0.10$, $s_0 = 0.20$ and a scale $\alpha_{\max} = 0.5$. The
shifts ensure that $a(z) \ge a_0 > 0$ and $s_k(z) \ge s_0 > 0$ at
initialization (avoiding the early-training collapse mode in which a learned
positive scale is initialized arbitrarily close to zero), and the bound on
$\alpha_k$ keeps the tanh perturbations moderate so that the linear term
$a(z)\,\log t$ dominates at the start of training and the flow does not
diverge before any gradient signal arrives. The intercept $b(z)$ and the knot
locations $\mu_k(z)$ are unconstrained.

\paragraph{Strict monotonicity.}
The Jacobian of~\eqref{eq:flow} with respect to $\log t$ is
\begin{equation}
    h'_\theta(\log t;\, z)
    \;\equiv\;
    \frac{\partial h_\theta}{\partial \log t}(\log t;\, z)
    \;=\;
    a(z) \;+\; \sum_{k=1}^{K} \frac{\alpha_k(z)}{s_k(z)}\,
        \operatorname{sech}^{2}\!\left(\frac{\log t - \mu_k(z)}{s_k(z)}\right).
    \label{eq:jacobian}
\end{equation}
Each $\operatorname{sech}^{2}$ term is non-negative, $\alpha_k(z), s_k(z) > 0$
by~\eqref{eq:reparam}, and $a(z) \ge a_0 > 0$, so
$h'_\theta(\log t;\, z) \ge a_0 > 0$ for all $\log t \in \mathbb{R}$ and all
$z$. Strict monotonicity is therefore enforced by construction (rather than
by a soft penalty), which guarantees that $h_\theta(\,\cdot\,;\, z)$ is a
diffeomorphism on $\mathbb{R}$ and that the change-of-variable formula below
holds without exception.

\paragraph{Closed-form density and survival.}
We take the base distribution to be standard normal, $U \sim \mathcal{N}(0, 1)$,
and define $T$ implicitly through $U = h_\theta(\log T;\, Z)$. Because
$h_\theta(\,\cdot\,; z)$ is strictly increasing,
$F_{T|Z}(t \mid z) = \Phi(h_\theta(\log t;\, z))$, and the chain rule
yields
\begin{align}
    \log f_{T|Z}(t \mid z)
    &\;=\;
    \log \phi\!\left(h_\theta(\log t;\, z)\right)
    \;+\;
    \log\, h'_\theta(\log t;\, z)
    \;-\;
    \log t,
    \label{eq:logpdf}\\[2pt]
    \log S_{T|Z}(t \mid z)
    &\;=\;
    \log \Phi\!\left(-\,h_\theta(\log t;\, z)\right),
    \label{eq:logsurv}
\end{align}
where $\phi$ and $\Phi$ are the standard-normal density and CDF. Both
expressions are evaluated entirely in the \emph{forward} direction
(equation~\eqref{eq:flow} plus its analytical derivative~\eqref{eq:jacobian}),
so the likelihood never invokes the inverse of the flow.

\paragraph{Numerical stability.}
Three implementation details prevent the well-known numerical failure modes
in the upper and lower tails of the likelihood:
\begin{itemize}
    \item In~\eqref{eq:logsurv} we evaluate $\log \Phi(-h)$ via
          \texttt{torch.special.ndtr} on the negated argument, which uses the
          symmetry $\Phi(-h) = 1 - \Phi(h)$ implicitly and avoids the
          catastrophic cancellation that would otherwise occur as
          $h \to +\infty$. The result is then clamped from below by
          $10^{-12}$ before the logarithm.
    \item In~\eqref{eq:logpdf} the Jacobian $h'_\theta$ is clamped from
          below by $10^{-12}$ before the logarithm. The clamp is never active
          at convergence because $h'_\theta \ge a_0 = 0.10$, but it makes
          the gradient well-defined under any pathological transient.
    \item All occurrences of $\log y$ first apply
          $\operatorname{clamp}(y, \ge 10^{-8})$ so that observation times
          arbitrarily close to zero do not produce $-\infty$.
\end{itemize}

\paragraph{Comparison with the parametric LogNormal alternative.}
Setting $K = 0$ in~\eqref{eq:flow} reduces $h_\theta$ to
$h_{\text{LN}}(\log t; z) = a(z)\,\log t + b(z)$, which is exactly the
conditional LogNormal model with location $-b(z)/a(z)$ and scale $1/a(z)$.
This is a strict special case of the proposed flow and is the "LogN-flow
ablation" reported in Section~\ref{sec:sim}: it isolates the contribution of
the $K$ tanh components ($K = 4$ in our experiments) under an otherwise
identical optimization pipeline.

\subsection{Censored joint log-likelihood}
\label{sec:mle}

Under conditionally independent censoring, the contribution of a single
observation $(x, \delta, z)$ to the joint log-likelihood factorizes as
\begin{align}
    \ell_{\theta_1, \theta_2}(x, \delta, z)
    &\;=\;
    \delta \,\Big[\log f_{T|Z}(x \mid z) + \log S_{C|Z}(x \mid z)\Big]
    \nonumber\\
    &\quad
    +\;
    (1 - \delta)\,\Big[\log S_{T|Z}(x \mid z) + \log f_{C|Z}(x \mid z)\Big].
    \label{eq:censmle}
\end{align}
Both flows enter additively, so on each forward pass we evaluate
$\log f_{T|Z}(x \mid z), \log S_{T|Z}(x \mid z)$ via flow $h_{\theta_1}$ and
$\log f_{C|Z}(x \mid z), \log S_{C|Z}(x \mid z)$ via flow $h_{\theta_2}$,
combine them according to~\eqref{eq:censmle}, and average over the batch:
\(
    \mathcal{L}_{\text{MLE}}(\theta_1, \theta_2)
    = -\,\frac{1}{n} \sum_{i=1}^{n} \ell_{\theta_1, \theta_2}(X_i, \delta_i, Z_i).
\)

\subsection{Soft--Nelson--Aalen Wasserstein regularizer}
\label{sec:reg}

The censored MLE in~\eqref{eq:censmle} is sample-level: it scores each
observation against its conditional density but does not directly constrain
the model-implied population survival curve to match the observed one. We
augment it with a population-level distributional constraint based on the
Nelson--Aalen estimator and its $1$-Wasserstein distance to a model-simulated
counterpart, made differentiable in $\theta$ by a soft-indicator
construction.

\paragraph{Time grid and observation-side anchor.}
Let $\tau_0 < \tau_1 < \cdots < \tau_M$ be a time grid spanning the empirical
$1$\textsuperscript{st}--$99$\textsuperscript{th} percentile of the observed
$X$, with $M = 80$ knots. We compute the Nelson--Aalen survival on this grid
from $\mathcal{D}$,
\[
    \widehat{S}_{\text{obs}}(\tau_m)
    \;=\;
    \exp\!\left(
        -\sum_{j \le m}
        \frac{
            \sum_{i} \mathbf{1}\{X_i \in [\tau_{j-1}, \tau_j)\}\, \delta_i
        }{
            \sum_{i} \mathbf{1}\{X_i \ge \tau_j\}
        }
    \right),
\]
which is held fixed throughout training (it is a function of the data only).

\paragraph{Soft simulated curve.}
For each batch and each $s = 1, \dots, S$ we draw fresh base noise
$U_T^{(s, i)}, U_C^{(s, i)} \sim \mathcal{N}(0, 1)$ for $i = 1, \dots, n$,
push it through the inverses of the two flows to obtain
$T_i^{(s)}, C_i^{(s)}$ (Section~\ref{sec:samp}), and clip both to the
training time interval $[t_{\min}, t_{\max}] = [10^{-4}, 10]$. The simulated
event time and a smooth event indicator are
\[
    X_i^{(s)} \;=\; \min\!\left(T_i^{(s)},\, C_i^{(s)}\right),
    \qquad
    \widetilde{\delta}_i^{(s)}
    \;=\;
    \sigma\!\left(\frac{C_i^{(s)} - T_i^{(s)}}{\tau_{\text{temp}}}\right),
\]
with $\sigma$ the logistic sigmoid and temperature
$\tau_{\text{temp}} = 0.25$. The corresponding soft Nelson--Aalen survival on
the same grid $\{\tau_m\}$ is
\[
    \widetilde{S}_{\text{sim}}^{(s)}(\tau_m)
    \;=\;
    \exp\!\left(
        -\sum_{j \le m}
        \frac{
            \sum_i \widetilde{\delta}_i^{(s)} \,
            \big[
                \sigma\!\left((X_i^{(s)} - \tau_{j-1}) / \alpha\right)
                - \sigma\!\left((X_i^{(s)} - \tau_j) / \alpha\right)
            \big]
        }{
            \sum_i \sigma\!\left((X_i^{(s)} - \tau_j) / \alpha\right)
            + \varepsilon
        }
    \right),
\]
where the soft at-risk count is a sigmoid step at temperature $\alpha = 0.20$
and the soft event count is a difference-of-sigmoids window at the same
temperature. The numerator stabilizer is $\varepsilon = 10^{-6}$.

\paragraph{Wasserstein objective.}
The $1$-Wasserstein distance between two survival curves on the real line
equals the absolute area between them,
\[
    W_1\!\left(\widehat{S}_{\text{obs}},\, \widetilde{S}_{\text{sim}}^{(s)}\right)
    \;=\;
    \int_{0}^{\infty}
        \big|
            \widehat{S}_{\text{obs}}(t) - \widetilde{S}_{\text{sim}}^{(s)}(t)
        \big|\,
    \mathrm{d}t,
\]
which we discretize on $\{\tau_m\}$ with the trapezoidal rule. Averaging over
$S$ simulations gives the regularizer
\(
    \mathcal{L}_{\text{NA-W}}(\theta_1, \theta_2)
    =
    \frac{1}{S} \sum_{s=1}^{S}
        W_1\!\left(\widehat{S}_{\text{obs}},\, \widetilde{S}_{\text{sim}}^{(s)}\right),
\)
and the total training objective is
\begin{equation}
    \mathcal{L}(\theta_1, \theta_2)
    \;=\;
    \mathcal{L}_{\text{MLE}}(\theta_1, \theta_2)
    \;+\;
    \lambda \cdot \mathcal{L}_{\text{NA-W}}(\theta_1, \theta_2),
    \label{eq:total}
\end{equation}
with $\lambda = 0.2$ and $S = 3$ in all experiments.

\subsection{Differentiable sampling via implicit reparameterization}
\label{sec:samp}

Computing $\mathcal{L}_{\text{NA-W}}$ requires backpropagating through the
samples
$T_i^{(s)} = h_{\theta_1}^{-1}(U_T^{(s, i)};\, Z_i)$ and similarly for $C$.
Because the inverse of~\eqref{eq:flow} has no closed form, we evaluate it
numerically and reattach the gradient via the
implicit-function construction of \citet{figurnov2018implicit}.

\paragraph{Inverse via vectorized bisection.}
Given a base draw $u \in \mathbb{R}$ and a covariate $z$, we seek
$\log t^{*}$ such that $h_\theta(\log t^{*};\, z) = u$. We initialize
$[\ell_0, h_0] = [\log t_{\min},\, \log t_{\max}]$ on the entire batch and
run $40$ steps of vectorized bisection,
\[
    m_k = \tfrac{1}{2}(\ell_k + h_k),
    \quad
    \ell_{k+1} = \begin{cases} m_k & h_\theta(m_k; z) < u \\ \ell_k & \text{else}\end{cases},
    \quad
    h_{k+1} = \begin{cases} h_k & h_\theta(m_k; z) < u \\ m_k & \text{else}\end{cases},
\]
returning $\log t^{*} = \tfrac{1}{2}(\ell_{40} + h_{40})$. The whole loop is
wrapped in \texttt{torch.no\_grad}, so it has zero gradient with respect to
$\theta$ and zero memory cost in the autograd graph. After $40$ iterations
the bracket width is $|h_{40} - \ell_{40}| \le |\log t_{\max} - \log t_{\min}|\, 2^{-40}
\approx 10^{-12}\,|\log t_{\max} - \log t_{\min}|$ in the worst case, far below
single-precision noise on $\log t^{*}$. This gives a numerically accurate
inverse but blocks gradient flow into the flow regularizer entirely.

\paragraph{Reattaching gradients.}
Because $h_\theta(\log t^{*};\, z) = u$ is an identity in $\theta$, the
implicit function theorem yields
$\partial \log t^{*} / \partial \theta
 = -(\partial h_\theta / \partial \theta) / h'_\theta$,
both evaluated at $\log t^{*}$. We realize this gradient via a one-line
substitution:
\begin{equation}
    \log t^{\diamond}
    \;=\;
    \operatorname{stop\_grad}(\log t^{*})
    \;-\;
    \frac{
        h_\theta\!\left(\log t^{*};\, z\right) - u
    }{
        \operatorname{stop\_grad}\!\left(h'_\theta(\log t^{*};\, z)\right)
    }.
    \label{eq:implicit}
\end{equation}
Numerically $\log t^{\diamond} = \log t^{*}$ to floating-point precision,
because $h_\theta(\log t^{*};\, z) = u$ at the bisection optimum. But the
autograd derivative of~\eqref{eq:implicit} with respect to $\theta$ equals
$-(\partial h_\theta / \partial \theta) / h'_\theta$ evaluated at
$\log t^{*}$, which is exactly the implicit-function gradient: the numerator
contributes $\partial h_\theta / \partial \theta$ through the (non-stopped)
forward call $h_\theta(\log t^{*};\, z)$, the denominator is stopped, and the
constant subtraction $\operatorname{stop\_grad}(\log t^{*})$ contributes
nothing to the gradient. Substituting $T = \exp(\log t^{\diamond})$ in the
regularizer therefore yields exact gradients without solving any system of
equations and without unrolling the bisection.

\paragraph{Training direction is forward, sampling direction is inverse.}
We emphasize the asymmetry: in the censored MLE
of~\eqref{eq:censmle}--\eqref{eq:logpdf}--\eqref{eq:logsurv}, the forward
direction is $\log t \mapsto u$, which is the closed-form direction of our
flow, so the likelihood needs only $h_\theta$ and its analytical Jacobian and
\emph{never} invokes the inverse. The inverse is invoked only in the
regularizer, where samples are generated by bisection and reattached to
gradients via~\eqref{eq:implicit}. This is in contrast to flow constructions
that compute the likelihood through the inverse direction with bisection,
where the derivative of the inverse must additionally be recovered (often
incorrectly) by automatic differentiation through the bisection itself; we
have observed this failure mode in a sibling implementation that motivated
the design here.

\subsection{Optimization}
\label{sec:opt}

We optimize~\eqref{eq:total} with full-batch Adam at learning rate $10^{-3}$
on the joint parameter set $\theta_1 \cup \theta_2$. After each backward pass
we apply gradient clipping at $\ell_2$-norm $5$ jointly to both flows; the
clip is rarely active at convergence but eliminates the rare exploding-gradient
events that occur in the first few epochs when the implicit gradient
in~\eqref{eq:implicit} is evaluated near a still-flat $h'_\theta$. We early-stop
on the training objective with patience $35$ and a relative tolerance of
$10^{-4}$, restoring the best-loss state at the end of training. Total
training time is approximately $20$ seconds per Monte Carlo replication on a
single CPU core for $n = 384$, $K = 4$, $S = 3$, and a maximum of $220$
epochs; convergence typically occurs around epoch $200$ and is uniform across
seeds.

\subsection{Hyperparameters}
\label{sec:hp}

For reproducibility, Table~\ref{tab:hp} lists every hyperparameter used in
the experiments of Section~\ref{sec:sim}. All values are shared across the
four data-generating processes and across all Monte Carlo replications;
nothing is tuned per scenario.

\begin{table}[t]
\caption{Hyperparameters for the conditional flow (\texttt{MonotoneFlow1D}).}
\label{tab:hp}
\centering
\small
\begin{tabular}{lll}
\toprule
Symbol & Value & Meaning \\
\midrule
$d$               & $1$ & covariate dimension in our simulations \\
hidden width      & $32$ & MLP $\psi_\theta$ hidden width (two layers) \\
$K$               & $4$ & number of tanh components in~\eqref{eq:flow} \\
$a_0$             & $0.10$ & shift in softplus reparam of $a(z)$ \\
$s_0$             & $0.20$ & shift in softplus reparam of $s_k(z)$ \\
$\alpha_{\max}$   & $0.5$ & scale on softplus reparam of $\alpha_k(z)$ \\
$[t_{\min}, t_{\max}]$ & $[10^{-4},\, 10]$ & training time clip \\
$M$               & $80$ & Nelson--Aalen grid size \\
$S$               & $3$ & MC simulations per batch in $\mathcal{L}_{\text{NA-W}}$ \\
$\tau_{\text{temp}}$ & $0.25$ & soft event-indicator temperature \\
$\alpha$          & $0.20$ & soft at-risk / event-window temperature \\
$\lambda$         & $0.2$ & weight of $\mathcal{L}_{\text{NA-W}}$ in~\eqref{eq:total} \\
optimizer         & Adam, lr $10^{-3}$ & full-batch \\
grad clip ($\ell_2$) & $5.0$ & joint over both flows \\
max epochs        & $220$ & with early stop \\
patience          & $35$ & early-stop patience \\
min$\Delta$       & $10^{-4}$ & early-stop relative tolerance \\
bisection iters   & $40$ & inverse evaluation in Section~\ref{sec:samp} \\
\bottomrule
\end{tabular}
\end{table}
